% !TeX root = ../sections/02_exercises.tex

\begin{exercise}
	有理数体上の多項式環 \(\mathbb{Q}[x]\) は整数環 \(\mathbb{Z}\) と同じく
	ユークリッド整域であり，従って PID，UFD である。
	では，整数環 \(\mathbb{Z}\) のかわりに有理数体上の多項式環
	\(\mathbb{Q}[x]\) で同じような問題を考えよう。
	
	中国式剰余定理，単因子の多項式版の問題。
	
	\begin{enumerate}
		\item[(A)]
		準同型写像
		\[
		\varphi:\mathbb{Q}[x]
		\longrightarrow
		\mathbb{Q}[x]/(x^2+x+1)
		\oplus
		\mathbb{Q}[x]/(x^2+2x+2)
		\]
		を
		\[
		\varphi(f(x))
		=
		\left(\widehat{f(x)},\overline{f(x)}\right)
		\]
		で定める。
		ただし，\(\widehat{f(x)}\) は \(f(x)\) を \(x^2+x+1\) で割った余り
		\((\)剰余類\()\) を表し，
		\(\overline{f(x)}\) は \(f(x)\) を \(x^2+2x+2\) で割った余り
		\((\)剰余類\()\) を表す。
		
		\begin{enumerate}
			\item
			\[
			\mathbb{Q}[x]/
			\left((x^2+x+1)(x^2+2x+2)\right)
			\simeq
			\mathbb{Q}[x]/(x^2+x+1)
			\oplus
			\mathbb{Q}[x]/(x^2+2x+2)
			\]
			を示せ。
			
			\item
			\(x^2+x+1\) で割った余りが \(x-1\) となり，
			\(x^2+2x+2\) で割った余りが \(x-2\) であるような
			有理数係数の多項式を求めよ。
		\end{enumerate}
		
		\item[(B)]
		行列
		\[
		A=
		\begin{pmatrix}
			x^2-1 & 2x+2\\
			x^2-x-2 & 3x+3
		\end{pmatrix}
		\]
		の単因子となる多項式 \(f_1(x),f_2(x)\) を求めよ。
		また，
		\[
		PAQ=\operatorname{diag}(f_1(x),f_2(x))
		\]
		が成り立つようなユニモジュラー行列 \(P,Q\) を求めよ。
	\end{enumerate}
\end{exercise}

\begin{background}
	この問題では，整数環 \(\mathbb{Z}\) で行っていた
	中国式剰余定理と単因子標準形の議論を，
	多項式環 \(\mathbb{Q}[x]\) 上で行う。
	
	\(\mathbb{Q}[x]\) はユークリッド整域である。
	したがって，\(\mathbb{Q}[x]\) は PID であり，
	任意のイデアルは一つの多項式によって生成される。
	また，ユークリッドの互除法を使うことができるため，
	整数の場合と同様に最大公約数やベズーの等式を扱うことができる。
	
	\[
	p(x)=x^2+x+1,
	\qquad
	q(x)=x^2+2x+2
	\]
	とおく。
	もし \(p(x)\) と \(q(x)\) が互いに素であれば，
	中国式剰余定理より
	\[
	\mathbb{Q}[x]/(p(x)q(x))
	\simeq
	\mathbb{Q}[x]/(p(x))
	\oplus
	\mathbb{Q}[x]/(q(x))
	\]
	が成り立つ。
	
	また，\(\mathbb{Q}[x]\) 上の行列についても，
	整数行列の場合と同様に単因子標準形を考えることができる。
	すなわち，行基本変形・列基本変形を用いて，
	\[
	PAQ=
	\operatorname{diag}(f_1(x),f_2(x))
	\]
	の形に変形する。
	ここで \(P,Q\) は \(\mathbb{Q}[x]\) 上のユニモジュラー行列であり，
	単因子は
	\[
	f_1(x)\mid f_2(x)
	\]
	を満たすように取る。
\end{background}

\begin{strategy}
	まず，
	\[
	p(x)=x^2+x+1,
	\qquad
	q(x)=x^2+2x+2
	\]
	とおく。
	
	\begin{enumerate}
		\item[(A1)]
		写像
		\[
		\varphi:\mathbb{Q}[x]\to
		\mathbb{Q}[x]/(p(x))
		\oplus
		\mathbb{Q}[x]/(q(x))
		\]
		の核を調べる。
		
		\[
		\varphi(f(x))=(0,0)
		\]
		であることは，
		\[
		p(x)\mid f(x),
		\qquad
		q(x)\mid f(x)
		\]
		であることと同値である。
		
		\(p(x)\) と \(q(x)\) が互いに素であることを示せば，
		\[
		p(x)q(x)\mid f(x)
		\]
		が従う。
		したがって，
		\[
		\Ker\varphi=(p(x)q(x))
		\]
		となる。
		
		その後，準同型定理を用いて
		\[
		\mathbb{Q}[x]/(p(x)q(x))
		\simeq
		\operatorname{Im}\varphi
		\]
		を得る。
		さらに中国式剰余定理により \(\varphi\) が全射であることを確認すれば，
		求める同型が従う。
		
		\item[(A2)]
		求めたい多項式を \(F(x)\) とする。
		条件は
		\[
		F(x)\equiv x-1 \pmod{p(x)},
		\qquad
		F(x)\equiv x-2 \pmod{q(x)}
		\]
		である。
		
		このような \(F(x)\) は，中国式剰余定理の構成を用いて求める。
		具体的には，ベズーの等式
		\[
		1=u(x)p(x)+v(x)q(x)
		\]
		を満たす \(u(x),v(x)\in\mathbb{Q}[x]\) を求める。
		
		このとき，
		\[
		v(x)q(x)\equiv 1 \pmod{p(x)},
		\qquad
		v(x)q(x)\equiv 0 \pmod{q(x)}
		\]
		であり，
		\[
		u(x)p(x)\equiv 0 \pmod{p(x)},
		\qquad
		u(x)p(x)\equiv 1 \pmod{q(x)}
		\]
		である。
		
		したがって，
		\[
		F(x)
		=
		(x-1)v(x)q(x)
		+
		(x-2)u(x)p(x)
		\]
		とおけば，指定された二つの剰余条件を同時に満たす。
		最後に必要なら \(p(x)q(x)\) で割った余りに直す。
		
		\item[(B)]
		行列
		\[
		A=
		\begin{pmatrix}
			x^2-1 & 2x+2\\
			x^2-x-2 & 3x+3
		\end{pmatrix}
		\]
		を \(\mathbb{Q}[x]\) 上の行列として考える。
		
		まず，各成分を因数分解する。
		特に，
		\[
		x^2-1=(x-1)(x+1),
		\qquad
		x^2-x-2=(x-2)(x+1),
		\]
		\[
		2x+2=2(x+1),
		\qquad
		3x+3=3(x+1)
		\]
		である。
		すべての成分に共通因子 \(x+1\) が現れることに注意する。
		
		単因子 \(f_1(x),f_2(x)\) を求めるには，
		\[
		\Delta_1=\gcd(\text{全成分})
		\]
		と
		\[
		\Delta_2=\gcd(\text{\(2\) 次小行列式})
		\]
		を計算する。
		\(2\times 2\) 行列の場合，\(\Delta_2\) は
		\[
		|\det A|
		\]
		に対応する。
		
		したがって，
		\[
		f_1(x)=\Delta_1,
		\qquad
		f_2(x)=\frac{\Delta_2}{\Delta_1}
		\]
		によって単因子を決定する。
		
		その後，実際に \(P,Q\) を求めるには，
		ユークリッドの互除法に対応する行基本変形・列基本変形を行い，
		まず左上に最大公約因子 \(f_1(x)\) を作る。
		次に，第1行・第1列の残りを消去し，
		右下に \(f_2(x)\) が残る形へ変形する。
	\end{enumerate}
\end{strategy}